\scanpage{253}

\mspage{129}

et de même

\[G=T_0.H=T'_0.H\qquad(\text{produits }1/2\text{ directs}),\tag*{(108)}\]

pour prouver (108), notons d'abord

\[T_0\cap H\subset T_G\cap H\deleted{\unsure{\subseteq}}\overset{(96)}{=}T_H\quad\text{donc}\]

\diagram{\begin{tikzcd}[ampersand replacement=\&,column sep=small]
{T_0\cap H=T_0\cap T_H=\operatorname{Ker}\bigl(} \&[-1.6em] T_H \arrow[r,hook,"\unsure{k|\kappa}"] \arrow[rr,bend right=22,swap,"\unsure{k|\kappa}"] \& T_G \arrow[r,"\unsure{\delta_0'|\delta_0}"] \& Z \&[-1.6em] {\bigr)=\{1\};}
\end{tikzcd}}

on voit de $=$ que $T'_0\cap H=\{1\}$. Pour voir que tout $g\in G$ s'écrit sous la forme $t_0.u$ ($t_0\in T_0,u\in H$), ou ce qui revient au $=$, sous la forme $u.t_0$ (puisque $T_0$ normalise $H$), OPS, comme $H$ est transitif sur $S$, quitte \deleted{\unsure{mod.}} \unsure{à} corriger \inserted{$g$} par un $u\in H$, que l'on a $gs_0=s_0$, donc on est ramené à la formule (107). Quitte à corriger par un élément de $T_0$ (puisque $T_0\xrightarrow{\delta|T_0} Z$ surjectif) OPS $\delta(g)=1$ i.e. $g\in G_0\cap B'_0$ et \deleted{comme on a vu que} quitte à corriger par un $z\in\unsure{\mathfrak{z}|\mathfrak{Z}}\subset T_0$, OPS $g\in H\cap B'_0\deleted{\unsure{=}}$, et il suffit de \deleted{voir} \inserted{noter} que

\[\begin{aligned}
H\cap B'_0&=U_0,\\
H\cap B'_1&=U_1,
\end{aligned}\tag*{(109)}\]

\deleted{\unsure{q}} ce qui achève de prouver (107) (108). On prouve de $=$

\[\left\{\begin{aligned}
B_{0G}&=T_0.B_{0H}\deleted{\unsure{=}}\\
B_{1G}&=T_{1}.B_{1H}
\end{aligned}\right.\ ,\tag*{(110)}\]

notons aussi

\[\left\{\begin{aligned}
\unsure{\hat B'_0|B'_0}\cap G_0&=\mathfrak{z}.U_0\\
B'_1\cap G_1&=\mathfrak{z}.U_1
\end{aligned}\right.\qquad\text{prod. }1/2\text{ directs}\tag*{(111)}\]

% archive sheet number: 252
